\chapter{Fase 3: elegir y medir los cinco lugares}
\label{cap:d-f3}
\textbf{Fecha}: 28 de septiembre de 2026. \textbf{Notas}: \codigo{01-regiones.md}, \codigo{05-planteamiento.md}.

\textbf{Qué se buscaba}: decidir qué lugares promueve la campaña con criterios medidos, y plantear el problema (variables,
actores y relaciones) con cifras.

\begin{opciones}[3.1 · De 8 regiones a 5 · 27 y 28-sep-2026]
\textbf{Primera versión (27-sep)}: Brandon pidió \emph{``mínimo 6 regiones''} sin los problemas de 2026, y eligió 8 destinos
(Chetumal, Calderitas–Oxtankah, la Ruta arqueológica, Maya Ka'an + Kantemó, Cobá + Punta Laguna, Muyil, Laguna
Milagros–Xul-Ha y Ribera del Río Hondo) más Cancún y Riviera Maya como emisoras.

\textbf{Revisión (28-sep)}: Brandon pidió no meter nada con sargazo, cierres o relación con Tulum, y enfocarse en 5.

\begin{center}
\small
\begin{tabularx}{\linewidth}{lY}
\encabezado \h{Región retirada} & \h{Motivo} \\
Cobá + Punta Laguna & Municipio de Tulum. Tenía buenos datos (74.4\,\% de capacidad probada sin usar), pero su imagen va
ligada a Tulum \\
\filagris Muyil & Cerrada del 4 de junio de 2024 al 10 de febrero de 2026 (INAH) y colinda con Tulum \\
Ribera del Río Hondo & Ninguna serie turística propia \\
\filagris Chacchoben & El 92\,\% de sus visitantes llega en crucero por Mahahual (franja de sargazo) \\
Kantunilkín–Chiquilá & Evidencia débil; es la puerta de Holbox \\
\bottomrule
\end{tabularx}
\normalsize
\end{center}

\textbf{Las 5}: Chetumal · Bahía Calderitas–Oxtankah · Ruta arqueológica del sur (Kohunlich, Dzibanché, Ichkabal) · Maya
Ka'an + Kantemó · Laguna Milagros–Xul-Ha. La quinta salió de las que quedaban: la Laguna no tiene sargazo ni cierres y está
en el mismo municipio que Chetumal.

\textbf{Evidencia (INAH)}: la zona de Tulum cayó 31.3\,\% en enero–julio de 2026 (de 692,946 a 476,247); la Ruta recibe
15.5 veces menos visitantes que Tulum (66,628 contra 1,031,443 en 2025).
\end{opciones}

\begin{opciones}[3.2 · El norte: ¿quitarlo o dejarlo como referencia? · 28-sep-2026]
\textbf{Elegida}: Cancún, Riviera Maya y Tulum \textbf{solo como referencia}, nunca como destino promovido.

\textbf{Descartado}: quitarlos por completo. \textbf{Por qué}: la única ocupación hotelera medida en 2025–2026 es la del
norte (DataTur), y sin ella no se puede medir de dónde se redistribuye ni comparar escenarios (criterios 6 y 9 de la
rúbrica).
\end{opciones}

\begin{opciones}[3.3 · La población: ¿por municipio o por localidad? · 28-sep-2026]
\textbf{Elegida}: \textbf{por localidad} (Censo 2020). \textbf{Por qué}: 4 de los 5 lugares están en Othón P. Blanco
(233,648 habitantes); con la cifra municipal tendrían el mismo denominador y no se podrían comparar.

\begin{center}
\small
\begin{tabularx}{\linewidth}{lYr}
\encabezado \h{Región} & \h{Localidades} & \h{Población} \\
Chetumal & Chetumal & 169,028 \\
\filagris Bahía & Calderitas & 5,551 \\
Laguna Milagros & Xul-Ha y Huay-Pix & 4,182 \\
\filagris Ruta arqueológica & Nicolás Bravo, Morocoy, Francisco Villa & 6,098 \\
Maya Ka'an + Kantemó & Felipe Carrillo Puerto, Tihosuco, Chunhuhub, Señor, Kantemó & 44,388 \\
\bottomrule
\end{tabularx}
\normalsize
\end{center}
\textbf{Riesgo declarado}: la lista es un criterio del proyecto, no una delimitación oficial. El código se detiene si una clave
no corresponde al nombre esperado (hay un ``Xul-Ha'' en Puerto Morelos).
\end{opciones}

\begin{opciones}[3.4 · El sur sin ocupación hotelera · 28-sep-2026]
\begin{center}
\small
\begin{tabularx}{\linewidth}{>{\raggedright\arraybackslash}p{0.32\linewidth}Y}
\encabezado \h{Opción} & \h{Por qué sí o por qué no} \\
\textbf{Medir la presión de llegada y mostrar ``sin dato oficial'' (elegida)} & Usa solo lo medido hasta 2026: Tren Maya
(3,843 bajaron en Chetumal en julio de 2026), cruces de Belice (49,097 en junio), INAH y cuartos \\
\filagris Estimar la ocupación con un modelo (\codigo{_est}) & Sería una cifra estimada donde nadie midió; con solo 36 meses de
Chetumal, además, poco confiable \\
\bottomrule
\end{tabularx}
\normalsize
\end{center}
\end{opciones}

\begin{opciones}[3.5 · El criterio de cierres · 28-sep-2026]
\textbf{La pregunta}: las zonas del sur también cerraron por obras en 2024 (Oxtankah de octubre de 2023 a octubre de 2024,
Kohunlich de marzo a diciembre de 2024, Dzibanché de febrero de 2024 a enero de 2025). Muyil salió por cierre: hacía falta
una regla igual para todas.

\textbf{Elegida}: una zona pasa si lleva \textbf{al menos 12 meses seguidos abierta} y no tuvo cierres en 2026. Pasan
Oxtankah (21 meses), Kohunlich (19) y Dzibanché (18); Muyil (6) no.

\textbf{Descartado}: ``ningún cierre desde 2024'', que habría sacado a la Ruta y a la Bahía y dejado solo 3 regiones.
\end{opciones}

\begin{opciones}[3.6 · Cómo se suman los cuartos del estado · 28-sep-2026]
\textbf{Hallazgo}: en SITUR-Q una zona no es la suma de sus miembros. En julio de 2026 la Riviera Maya tenía 59,632 cuartos,
contra 22,923 de Playa del Carmen + Tulum: \textbf{36,709 de diferencia} (lugares que no se publican por separado).

\textbf{Elegida}: el total estatal suma destinos + zonas (\textbf{140,664 cuartos}). \textbf{Descartado}: destinos +
miembros, que perdería esos 36,709 cuartos.
\end{opciones}

\section*{Resultado: el problema medido}
\textbf{La tabla de criterios} (sargazo, cierres, ocupación, presión por residente, oferta, servicios y series disponibles):
las 5 regiones pasan. Donde hay dato, sus hoteles se llenaron mucho menos en 2024 (Chetumal 58\,\%, Maya Ka'an 39\,\%) que
los del norte (74–77\,\%).

\textbf{La relación central}, con cuotas y el índice de Herfindahl-Hirschman:
\begin{center}
\small
\begin{tabular}{lrrr}
\encabezado \h{Dimensión} & \h{Concentración} & \h{Parte de los 5} & \h{Contra su población} \\
Llegadas en avión (2024) & 0.811 & 1.4\,\% & 0.11 \\
Cuartos de hotel (jul-2026) & 0.219 & 1.7\,\% & 0.14 \\
Visitantes INAH (2025) & 0.276 & 4.1\,\% & 0.33 \\
Negocios turísticos & 0.133 & 14.4\,\% & 1.17 \\
Población (2020) & 0.225 & 12.3\,\% & 1.00 \\
\end{tabular}
\normalsize
\end{center}
\textbf{Lectura}: los negocios siguen a la población; los visitantes, no. Es capacidad ociosa.

\begin{opciones}[Un hallazgo que no se escondió]
Kohunlich (+13.5\,\%) y Dzibanché (+69.2\,\%) crecen en 2026, pero la Ruta completa bajó 8.2\,\% porque Ichkabal cayó
27.5\,\%. \textbf{Decisión}: no destacar en la portada solo las zonas que suben, porque sería escoger datos a conveniencia.
\end{opciones}
